\documentclass[../../main]{subfiles}

\begin{document}

\section{Operações sobre Relações}
\label{sec:matematica:relacoes:operacoes-sobre-relacoes}

\begin{defn}[União de Relações]
	\label{def:matematica:relacoes:operacoes:uniao}

	Sejam \(R\) e \(S\) relações sobre um conjunto \(A\).

	Define-se a união de \(R\) e \(S\) por

	\[
		R \cup S
		=
		\{
		(a,b)
		\mid
		(a,b)\in R
		\lor
		(a,b)\in S
		\}.
	\]

\end{defn}

\begin{defn}[Interseção de Relações]
	\label{def:matematica:relacoes:operacoes:intersecao}

	Sejam \(R\) e \(S\) relações sobre um conjunto \(A\).

	Define-se a interseção de \(R\) e \(S\) por

	\[
		R \cap S
		=
		\{
		(a,b)
		\mid
		(a,b)\in R
		\land
		(a,b)\in S
		\}.
	\]

\end{defn}

\begin{defn}[Diferença de Relações]
	\label{def:matematica:relacoes:operacoes:diferenca}

	Sejam \(R\) e \(S\) relações sobre um conjunto \(A\).

	Define-se a diferença entre \(R\) e \(S\) por

	\[
		R \setminus S
		=
		\{
		(a,b)
		\mid
		(a,b)\in R
		\land
		(a,b)\notin S
		\}.
	\]

\end{defn}

\begin{defn}[Relação Inversa]
	\label{def:matematica:relacoes:operacoes:inversa}

	Seja \(R\) uma relação de \(A\) em \(B\).

	A relação inversa de \(R\), denotada por

	\[
		R^{-1},
	\]

	é definida por

	\[
		R^{-1}
		=
		\{
		(b,a)
		\mid
		(a,b)\in R
		\}.
	\]

\end{defn}

\begin{ex}

	Se

	\[
		R
		=
		\{
		(1,2),
		(1,3),
		(2,3)
		\},
	\]

	então

	\[
		R^{-1}
		=
		\{
		(2,1),
		(3,1),
		(3,2)
		\}.
	\]

\end{ex}

\begin{prop}
	\label{prop:matematica:relacoes:operacoes:dupla-inversa}

	Para toda relação \(R\),

	\[
		(R^{-1})^{-1}
		=
		R.
	\]

\end{prop}

\begin{proof}

	Seja

	\[
		(a,b)\in (R^{-1})^{-1}.
	\]

	Pela definição de relação inversa,

	\[
		(b,a)\in R^{-1}.
	\]

	Novamente pela definição,

	\[
		(a,b)\in R.
	\]

	A inclusão recíproca é análoga.

\end{proof}

\begin{prop}
	\label{prop:matematica:relacoes:operacoes:inversa-uniao}

	Para quaisquer relações \(R\) e \(S\),

	\[
		(R\cup S)^{-1}
		=
		R^{-1}\cup S^{-1}.
	\]

\end{prop}

\begin{proof}

	Segue diretamente das definições de união e inversa.

\end{proof}

\begin{prop}
	\label{prop:matematica:relacoes:operacoes:inversa-intersecao}

	Para quaisquer relações \(R\) e \(S\),

	\[
		(R\cap S)^{-1}
		=
		R^{-1}\cap S^{-1}.
	\]

\end{prop}

\begin{proof}

	Segue diretamente das definições.

\end{proof}

\subsection{Composição de Relações}

\begin{defn}[Composição de Relações]
	\label{def:matematica:relacoes:operacoes:composicao}

	Sejam

	\[
		R \subseteq A\times B
	\]

	e

	\[
		S \subseteq B\times C.
	\]

	A composição de \(R\) com \(S\), denotada por

	\[
		S\circ R,
	\]

	é a relação definida por

	\[
		S\circ R
		=
		\{
		(a,c)
		\mid
		(\exists b\in B)
		\bigl(
		(a,b)\in R
		\land
		(b,c)\in S
		\bigr)
		\}.
	\]

\end{defn}

\begin{obs}

	A composição relaciona \(a\) com \(c\) sempre que existir um elemento
	intermediário \(b\) ligando ambas as relações. :contentReference[oaicite:1]{index=1}

\end{obs}

\begin{ex}

	Considere

	\[
		R
		=
		\{
		(1,2),
		(2,3)
		\}
	\]

	e

	\[
		S
		=
		\{
		(2,a),
		(3,b)
		\}.
	\]

	Então

	\[
		S\circ R
		=
		\{
		(1,a),
		(2,b)
		\}.
	\]

\end{ex}

\begin{prop}[Associatividade da Composição]
	\label{prop:matematica:relacoes:operacoes:associatividade-composicao}

	Sejam \(R\), \(S\) e \(T\) relações compatíveis.

	Então

	\[
		T\circ(S\circ R)
		=
		(T\circ S)\circ R.
	\]

\end{prop}

\begin{proof}

	Seja

	\[
		(a,d)\in T\circ(S\circ R).
	\]

	Então existem \(b\) e \(c\) tais que

	\[
		(a,b)\in R,
	\]

	\[
		(b,c)\in S,
	\]

	e

	\[
		(c,d)\in T.
	\]

	Essas mesmas condições caracterizam

	\[
		(a,d)\in (T\circ S)\circ R.
	\]

	Portanto as relações são iguais. :contentReference[oaicite:2]{index=2}

\end{proof}

\begin{prop}
	\label{prop:matematica:relacoes:operacoes:inversa-composicao}

	Para relações compatíveis \(R\) e \(S\),

	\[
		(S\circ R)^{-1}
		=
		R^{-1}\circ S^{-1}.
	\]

\end{prop}

\begin{proof}

	Seja

	\[
		(c,a)\in (S\circ R)^{-1}.
	\]

	Então

	\[
		(a,c)\in S\circ R.
	\]

	Existe \(b\) tal que

	\[
		(a,b)\in R
		\quad\text{e}\quad
		(b,c)\in S.
	\]

	Logo

	\[
		(b,a)\in R^{-1}
		\quad\text{e}\quad
		(c,b)\in S^{-1}.
	\]

	Portanto

	\[
		(c,a)\in R^{-1}\circ S^{-1}.
	\]

	A inclusão recíproca é análoga.

\end{proof}

\begin{defn}[Potência de uma Relação]
	\label{def:matematica:relacoes:operacoes:potencia}

	Seja \(R\) uma relação sobre um conjunto \(A\).

	Define-se

	\[
		R^1 = R
	\]

	e, para todo \(n\ge 1\),

	\[
		R^{n+1}
		=
		R^n \circ R.
	\]

\end{defn}

\begin{ex}

	Se

	\[
		R
		=
		\{
		(1,2),
		(2,3)
		\},
	\]

	então

	\[
		R^2
		=
		\{
		(1,3)
		\}.
	\]

\end{ex}

\end{document}
